\title{Noesis: A Protocol for Trading and Fulfilling Machine Intelligence as a Commodity}

\author{Giacinto Paolo Saggese$^{*}$}
\thanks{$^{*}$ The University of Maryland, College Park, MD 20742, USA, gsaggese@umd.edu}

\date{\today}

\begin{abstract}
Applications and autonomous agents increasingly consume large language model
(LLM) inference as a raw input, in the same way that factories consume
electricity. Today this input is purchased through static, per-model price
lists set unilaterally by a small number of providers, with no price
discovery, no verification that the delivered service matched what was paid
for, and no mechanism to route a request to only as much capability as it
actually requires.
%
We propose \Noesis, a protocol that treats machine
intelligence as a fungible commodity and organizes its exchange around two
cooperating components.
% TODO(gp): Market -> Exchange?
\NoesisMarket{} is a periodic call auction that
matches buyers, who bid for a bundle of task volume, capability tier,
latency, and reliability, against sellers, who ask a price for delivering
that bundle, clearing at a single uniform price per capability tier.
%
\NoesisServer{} is the fulfillment and monitoring layer: an API gateway that
executes the matched contracts against real model providers, meters whether
the delivered capability, latency, and reliability met the contracted terms,
and feeds the result back into the market as a reputation and pricing signal.
As a byproduct, \NoesisServer{} logs every request and estimates how much
capability each request actually needs, enabling difficulty-aware routing and
answer fusion that trade cost for accuracy in opposite directions, and
enabling distillation of the collected data into smaller, cheaper models.
% TODO(ai_gp): Do not talk about distillation
% TODO(ai_gp): Make it fully anonymized
We formalize the contract, bid, and ask notation; state the market-clearing
problem as a linear program and establish that a clearing price exists under
standard monotonicity assumptions; and describe the request-level fulfillment
statistics needed to attribute a shortfall to a specific seller rather than
to noise. We compare \Noesis{} against direct provider APIs, pass-through
gateways, and raw-compute spot markets, and lay out the open mechanism-design
questions and a phased implementation plan for a prototype.
\end{abstract}

\maketitle

\setcounter{tocdepth}{2}
\tableofcontents

% ##############################################################################
\section{Introduction}
\label{sec:introduction}

\Noesis{} is a protocol for exchanging machine intelligence, understood as LLM
inference capacity delivered at a stated capability, latency, and reliability
level, between buyers who need it and sellers who supply it. It brings together
price discovery through an open auction, a contract that bundles quality
attributes rather than a single price-per-token, and a fulfillment layer that
measures and reports whether the contract was actually honored.

% ==============================================================================
\subsection{Machine intelligence as a commodity}
\label{sec:intelligence_as_commodity}

Electricity, natural gas, and bandwidth all share a property that makes them
tradeable as commodities: any unit supplied by any seller is, for practical
purposes, interchangeable with any unit supplied by any other seller, once
quality attributes (voltage and frequency for electricity, calorific value for
gas, latency and jitter for bandwidth) are specified and metered. LLM inference
exhibits a version of the same property. A request answered at a stated
capability tier, under a stated latency bound, with a stated reliability, is,
from the buyer's perspective, largely interchangeable across providers that can
meet those terms, even though the underlying models, weights, and hardware
differ completely.

Unlike electricity, however, the traded unit cannot be reduced to a single
scalar such as kilowatt-hours. Raw compute, measured in FLOPs or GPU-hours,
is observable but is not what a buyer wants to purchase: two providers can
spend the same compute on the same request and produce answers of very
different quality, and a buyer typically does not know, or care, how much
compute was spent to produce a given answer. The good that a buyer actually
wants to buy is better described as a bundle: a capability level (how good
the answer is expected to be), a latency bound (how long the answer takes),
and a reliability level (how often the first two conditions are actually
met). Section~\ref{sec:bundling} returns to this point, which drives most
of the protocol's design.

% ==============================================================================
\subsection{The Noesis protocol}
\label{sec:protocol_overview}

\Noesis{} is organized around two cooperating components.

\begin{itemize}
  \item \textbf{\NoesisMarket}: a periodic call auction that matches buyers
    and sellers into contracts. Every $T$ minutes, buyers submit bids and
    sellers submit asks, bucketed by capability tier; within each tier, bids
    and asks are matched at a single uniform clearing price, subject to
    latency and reliability compatibility (Section~\ref{sec:noesis_market}).

  \item \textbf{\NoesisServer}: the fulfillment and monitoring layer. It is
    an LLM API gateway, modeled on multi-provider routers such as
    OpenRouter\footnote{\url{https://openrouter.ai}}, that executes the
    requests matched by a \NoesisMarket{} contract against the underlying
    providers, logs every prompt/response pair, and measures whether the
    delivered capability, latency, and reliability met the contract terms.
    Violations are reported back to \NoesisMarket{} as a reputation and
    pricing signal (Section~\ref{sec:noesis_server}).
\end{itemize}

The two components close a loop that is missing from today's LLM API
market: a contract is not merely sold, it is metered, and the metering
result changes the terms under which the same seller can sell in the
future. \NoesisServer{} additionally reuses the logged traffic for two
purposes orthogonal to fulfillment: difficulty-aware routing, which sends a
request to the cheapest provider expected to satisfy it, and answer fusion,
which combines several providers' answers when higher reliability is worth
the extra cost.

% ==============================================================================
\subsection{Contributions}
\label{sec:contributions}

The innovative features of \Noesis{} are:

\begin{enumerate}
  \item \textbf{Price discovery through auction, not price lists}:
    \NoesisMarket{} clears buyers and sellers through a periodic call
    auction instead of relying on prices unilaterally posted by a small
    number of providers.

  \item \textbf{A bundled, verifiable good}: the traded unit is a
    (capability, latency, reliability) bundle rather than a raw compute
    quantity, so the market quotes and clears on ``how good'', not only
    ``how much''.

  \item \textbf{Verified fulfillment}: \NoesisServer{} meters every
    request executed under a contract and computes a measured reliability
    and latency, so a contract's terms are checked against evidence rather
    than taken on faith.

  \item \textbf{Reputation and pricing feedback}: a seller whose measured
    performance falls short of its ask is flagged, and the resulting
    reputation signal feeds back into eligibility and pricing in future
    auction rounds, in the spirit of non-performance penalties in capacity
    markets.

  \item \textbf{Difficulty-aware routing}: \NoesisServer{} estimates how
    much capability a request actually needs and routes accordingly,
    spending intelligence only where it is needed.

  \item \textbf{Answer fusion as a complementary mode}: for requests where
    reliability matters more than cost, \NoesisServer{} can query several
    providers and combine their answers, trading cost for accuracy in the
    direction opposite to routing.

  \item \textbf{A reusable byproduct dataset}: every logged
    prompt/response pair is available for evaluation and distillation,
    without the market or the buyer having to do anything beyond normal
    use of the protocol.

  \item \textbf{Provider-agnostic liquidity pooling}: many sellers compete
    for the same demand behind a single interface, instead of buyers being
    locked into a single provider's price list.

  \item \textbf{Tiered commodity pricing}: capability tiers clear
    independently in each auction round, analogous to peak and off-peak
    pricing in electricity markets.
\end{enumerate}

% ==============================================================================
\subsection{Comparison with alternative solutions}
\label{sec:comparison_table}

Table~\ref{tab:comparison} compares \Noesis{} against three alternative ways
of obtaining LLM inference: calling a provider's API directly, using a
pass-through multi-provider gateway (e.g., OpenRouter-style routers), and
buying raw compute on a cloud spot market (e.g., spot GPU instances). The
comparison is discussed in detail in Sections~\ref{sec:market_design}
through \ref{sec:noesis_server}.

% TODO(ai_gp): Make the first row and column bold
% TODO(ai_gp): Make the symbols more readable by using darker colors
\begin{samepage}
  \begin{table}[H]
    \centering
    \begin{tabular}{l|c|c|c|c|}
                                              & \emph{Direct API} & \emph{Gateway} & \emph{Spot compute} & \emph{\Noesis} \\
      \hline
      Price discovery via open auction       & \nbad             & \nbad          & \ngood               & \ngood         \\
      \hline
      Bundled quality guarantee              & \npartial         & \npartial      & \nbad                & \ngood         \\
      \hline
      Verified fulfillment / SLA monitoring  & \nbad             & \nbad          & \npartial             & \ngood         \\
      \hline
      Reputation-based accountability        & \nbad             & \nbad          & \npartial             & \ngood         \\
      \hline
      Difficulty-aware cost routing          & \nbad             & \npartial      & \nbad                 & \ngood         \\
      \hline
      Multi-provider liquidity pooling       & \nbad             & \ngood         & \npartial             & \ngood         \\
      \hline
      Reusable evaluation/training dataset   & \npartial         & \npartial      & \nbad                 & \ngood         \\
      \hline
      Single-API ease of integration         & \ngood            & \ngood         & \nbad                 & \ngood         \\
      \hline
    \end{tabular}
    \caption{Comparison of \Noesis{} against alternative ways of obtaining
      LLM inference.}
    \label{tab:comparison}
  \end{table}
\end{samepage}

Cloud spot markets already apply auction-based price discovery, but to raw
compute rather than to a capability/latency/reliability bundle, so they
score poorly on the bundled-quality and difficulty-routing rows: a spot GPU
market has no notion of ``how good'' the resulting answer is. Gateways pool
liquidity across providers and are easy to integrate, but set prices
unilaterally rather than clearing an auction, and typically do not meter
fulfillment against a contract. \Noesis{} is designed to combine the
liquidity pooling and ease-of-integration of a gateway with the price
discovery of an auction market and the verifiability of a metered service
contract.
